%
\hsection{Integer Arithmetic / Rechnen mit Ganzzahlen}%
%
\begin{question}{Data type(s) / Datentyp(en)}%
\qENDE{%
Which data type(s) does \python~3 provide for whole numbers, i.e., integer numbers?\pythonIdx{int]}%
}{%
Welchen Datentyp bzw.\ welche Datentypen stellt \python~3 für Ganzzahlen zur Verfügung?%
}%
\end{question}%
%
\begin{question}{Range / Wertebereich}%
\qENDE{%
What are the theoretical \emph{and} practical limits for the largest integer number that can be represented with \python~3?%
}{%
Was sind die theoretischen \emph{und} praktischen Grenzen für die größte Ganzzahl, die mit \python~3 dargestellt werden kann?%
}%
\end{question}%
%
\begin{question}{Division}%
\qENDE{%
Explain the difference between the \pythonilIdx{//} and \pythonilIdx{/} operators in \python~3.%
}{%
Was ist der Unterschied zwischen den Operatoren \pythonilIdx{//} und \pythonilIdx{/} in \python~3?%
}%
\end{question}%
%
\begin{question}{Remainder / Rest}%
\qENDE{%
How can we compute the remainder of the division of two integer variables \pythonil{a} and \pythonil{b} in \python?
}{%
Wie können wir den Rest der Division zweier Ganzzahl-Variablen \pythonil{a} und \pythonil{b} in \python\ berechnen?%
}%
\end{question}%
%
\begin{question}{**-Operator}%
\qENDE{%
What is are the results of \pythonil{3 ** 2} and \pythonil{2 ** 4} in \python?\pythonIdx{**}%
}{%
Was sind die Ergebnisse von \pythonil{3 ** 2} und \pythonil{2 ** 4} in \python?\pythonIdx{**}%
}%
\end{question}%
%
\begin{question}{Program Understanding / Programm Verstehen}%
\gitLoadPython{int_02}{}{questions/int_02.py}{}%
\listingPython{int_02}{The program \programUrl{int_02}}%
\qENDE{%
Which output will the program \programUrl{int_02} in \cref{lst:int_02} produce?%
}{%
Welche Ausgabe produziert das Programm \programUrl{int_02} in \cref{lst:int_02}?%
}%
\end{question}%
%
\begin{question}{Program Understanding / Programm Verstehen}%
\gitLoadPython{int_01}{}{questions/int_01.py}{}%
\listingPython{int_01}{The program \programUrl{int_01}}%
%
\qENDE{%
\begin{enumerate}%
\item Explain what the program \programUrl{int_01} in \cref{lst:int_01} is \emph{supposed} to do.%
\item Does it achieve this? If not, why?%
\item What will happen if we execute this program?%
\item Find at least two errors in this program.%
\end{enumerate}%
%
}{%
\begin{enumerate}%
\item Erklären Sie, was das Programm \programUrl{int_01} in \cref{lst:int_01} machen \emph{soll}.%
\item Macht es das auch? Wenn nicht, warum?
\item Was wird passieren, wenn wir dieses Programm ausführen?%
\item Finden Sie mindestens zwei Fehler in diesem Programm.%
\end{enumerate}%
}%
\end{question}%
%
\begin{question}{Calculation / Berechnung}%
\qENDE{%
Write a \python~3 program computing the integer value~$Z=3+9^7-\frac{8}{2}$ using \emph{only} integer arithmetics. %
Notice that all the operations must actually be performed by the program. %
You cannot just compute the result on paper and just print that result\dots%
}{%
Schreiben Sie ein \python~3-Programm, das den Ganzzahlwert~$Z=3+9^7-\frac{8}{2}$ berechnet und dabei \emph{nur} Ganzzahlarithmetik verwendet. %
Beachten Sie, dass alle Rechenoperationen wirklich vom Programm ausgeführt werden müssen. %
Sie dürfen die Formel nicht einfach auf dem Papier ausrechnen und nur das Ergebnis ausgeben\dots%
}%
\end{question}%
%
\endhsection%
%
